\documentclass[11pt, a4paper]{article}

% bfw-style.tex — gemeinsame Style-Definitionen für alle Handouts/Aufgabenhefte
% Voraussetzung: Kompilation mit xelatex (wegen fontspec)

\usepackage[ngerman]{babel}
% Babel-ngerman macht " zu einem aktiven Shorthand (z.B. für "a -> ä). In unseren
% Texten verwenden wir " als normales Zeichen — daher Shorthand komplett ausschalten.
\shorthandoff{"}
\usepackage{geometry}
\geometry{a4paper, top=2.4cm, bottom=2.4cm, left=2.3cm, right=2.3cm,
          headheight=15pt, headsep=0.7cm, footskip=1.1cm}

\usepackage{fontspec}
% Fallback-Kette: Source Sans Pro -> Calibri -> Segoe UI -> Latin Modern Sans
% (Latin Modern ist immer mit MiKTeX vorhanden)
\IfFontExistsTF{Source Sans Pro}{%
  \setmainfont{Source Sans Pro}[Ligatures=TeX]%
}{\IfFontExistsTF{Calibri}{%
  \setmainfont{Calibri}[Ligatures=TeX]%
}{\IfFontExistsTF{Segoe UI}{%
  \setmainfont{Segoe UI}[Ligatures=TeX]%
}{%
  \setmainfont{Latin Modern Sans}[Ligatures=TeX]%
}}}
\IfFontExistsTF{Source Code Pro}{%
  \setmonofont{Source Code Pro}[Scale=0.92]%
}{\IfFontExistsTF{Consolas}{%
  \setmonofont{Consolas}[Scale=0.92]%
}{%
  \setmonofont{Latin Modern Mono}[Scale=0.92]%
}}

% Gesperrte Versalien für Labels/Titelbalken (Kapitälchen-Ersatz, funktioniert
% mit jeder OpenType-Schrift — Latin Modern Sans hat keine echten Kapitälchen)
\newcommand{\bfwlabelfont}{\footnotesize\bfseries\addfontfeatures{LetterSpace=8.0}}

\usepackage{xcolor}
% Farbpalette: hell, freundlich, modern
\definecolor{bfwPrimary}{HTML}{0EA5A4}    % Petrol/Türkis - Primär
\definecolor{bfwPrimaryDark}{HTML}{0F766E} % dunkler Petrol für Headings
\definecolor{bfwAccent}{HTML}{F59E0B}     % Amber - Akzent
\definecolor{bfwSuccess}{HTML}{10B981}    % Grün - Erfolg / Lösung
\definecolor{bfwWarn}{HTML}{F97316}       % Orange - Achtung
\definecolor{bfwInk}{HTML}{0F172A}        % fast Schwarz
\definecolor{bfwInkSoft}{HTML}{475569}    % gedämpftes Grau
\definecolor{bfwBgSoft}{HTML}{F8FAFC}     % off-white
\definecolor{bfwBgTeal}{HTML}{ECFEFF}     % helles Türkis
\definecolor{bfwBgAmber}{HTML}{FEF3C7}    % helles Gelb
\definecolor{bfwBgRose}{HTML}{FFE4E6}     % helles Rosé für Eselsbrücken
\definecolor{bfwTabHead}{HTML}{E4F3F2}    % zarte Kopfzeilen-Hinterlegung für Tabellen

\usepackage[colorlinks=true, linkcolor=bfwPrimaryDark, urlcolor=bfwPrimaryDark]{hyperref}
\usepackage{lastpage}
\usepackage{enumitem}
\setlist[itemize]{leftmargin=*, itemsep=2pt, topsep=2pt}
\setlist[enumerate]{leftmargin=*, itemsep=2pt, topsep=2pt}

\usepackage[most]{tcolorbox}
\usepackage{booktabs}
\usepackage{colortbl}
\usepackage{tikz}
\usetikzlibrary{decorations.pathmorphing, positioning, shapes.geometric, arrows.meta}

% ---------------------------------------------------------------------------
% Einheitliches Tabellen-Design
%   - etwas mehr Zeilenluft, dezente graue Linien (wirkt auch mit \hline ruhiger)
%   - Kopfzeile einer Tabelle bei Bedarf zart hinterlegen: \rowcolor{bfwTabHead}
% ---------------------------------------------------------------------------
\renewcommand{\arraystretch}{1.25}
\arrayrulecolor{bfwInkSoft!50}
\setlength{\heavyrulewidth}{1pt}
\setlength{\lightrulewidth}{0.5pt}

% ---------------------------------------------------------------------------
% Boxen — klare farbige Titelbalken mit gesperrten Versal-Labels
% (Titel bleibt per Option überschreibbar: \begin{merksatz}[title={...}])
% ---------------------------------------------------------------------------
\tcbset{bfwbox/.style={%
  enhanced, breakable,
  boxrule=0.5pt, arc=2.5pt,
  left=10pt, right=10pt, top=8pt, bottom=8pt,
  toptitle=4pt, bottomtitle=4pt,
  titlerule=0pt,
  fonttitle=\bfwlabelfont,
  coltitle=white
}}

% Merkkästchen — wichtige Definition / Kernpunkt
\newtcolorbox{merksatz}[1][]{%
  bfwbox,
  colback=bfwBgTeal,
  colframe=bfwPrimaryDark,
  colbacktitle=bfwPrimaryDark,
  title={MERKE},
  #1
}

% Eselsbrücke — Mnemoniks
\newtcolorbox{eselsbruecke}[1][]{%
  bfwbox,
  colback=bfwBgRose,
  colframe=bfwAccent!85!black,
  colbacktitle=bfwAccent!85!black,
  title={ESELSBRÜCKE},
  #1
}

% Beispiel-Box
\newtcolorbox{beispiel}[1][]{%
  bfwbox,
  colback=bfwBgSoft,
  colframe=bfwInkSoft,
  colbacktitle=bfwInkSoft,
  title={BEISPIEL},
  #1
}

% Lösungs-Box (für Aufgabenhefte)
\newtcolorbox{loesung}[1][]{%
  bfwbox,
  colback=bfwSuccess!7,
  colframe=bfwSuccess!65!black,
  colbacktitle=bfwSuccess!65!black,
  title={LÖSUNG},
  #1
}

% Achtung-Box — typische Fehler / Stolperfallen
\newtcolorbox{achtung}[1][]{%
  bfwbox,
  colback=bfwWarn!6,
  colframe=bfwWarn!85!black,
  colbacktitle=bfwWarn!85!black,
  title={ACHTUNG},
  #1
}

% Formel-Box — für zentrale Formeln (groß, ruhig, ohne Titelbalken)
\newtcolorbox{formelbox}[1][]{%
  enhanced,
  colback=bfwBgTeal,
  colframe=bfwPrimaryDark,
  boxrule=1pt, arc=3pt,
  left=10pt, right=10pt, top=6pt, bottom=6pt,
  #1
}

% Glossar-Eintrag
\newcommand{\glossareintrag}[2]{%
  \par\noindent\textbf{\color{bfwPrimaryDark}#1} \enspace #2\par\smallskip
}

% ---------------------------------------------------------------------------
% Abschnitts-Design: farbiges Nummern-Quadrat + feine Linie unter \section,
% dezent farbige \subsection/\subsubsection
% ---------------------------------------------------------------------------
\usepackage{titlesec}
\newcommand{\bfwSecBadge}[1]{%
  \begin{tikzpicture}[baseline={([yshift=-0.62em]n.center)}]
    \node[fill=bfwPrimaryDark, text=white, font=\normalsize\bfseries,
          minimum width=1.55em, minimum height=1.55em, inner sep=1pt,
          rounded corners=1.5pt] (n) {#1};
  \end{tikzpicture}}
\titleformat{\section}
  {\Large\bfseries\color{bfwPrimaryDark}}
  {\bfwSecBadge{\thesection}}{0.6em}{}
  [\vspace{-0.2em}{\color{bfwPrimary!60}\titlerule[0.8pt]}]
\titleformat{\subsection}
  {\large\bfseries\color{bfwPrimaryDark}}
  {\textcolor{bfwPrimary}{\thesubsection}}{0.5em}{}
\titleformat{\subsubsection}
  {\normalsize\bfseries\color{bfwInk}}
  {\textcolor{bfwPrimary}{\thesubsubsection}}{0.4em}{}
\titlespacing*{\section}{0pt}{1.6em}{1em}
\titlespacing*{\subsection}{0pt}{1.2em}{0.5em}

% TikZ-Stil "skizzenhaft" — etwas weicher als Rough.js, aber stabil
\tikzset{
  roughbox/.style = {
    draw=bfwPrimaryDark, thick, fill=bfwBgTeal,
    rounded corners=4pt, inner sep=6pt
  },
  roughaccent/.style = {
    draw=bfwAccent, thick, fill=bfwBgAmber,
    rounded corners=4pt, inner sep=6pt
  },
  roughline/.style = {
    draw=bfwInkSoft, thick, -{Stealth[length=2.5mm]}
  }
}

% Bullet-Symbole (bewusst kein FontAwesome — vermeidet Font-Installations-Probleme)
\newcommand{\faLightbulb}{$\bullet$}
\newcommand{\faBookmark}{$\bullet$}

% ---------------------------------------------------------------------------
% Kopf-/Fußzeile
%   Kopf links:  Wortlogo „Wirtschaftsfachwirt"
%   Kopf rechts: Dokument-Kurztext, gesetzt via \kopfzeile{...}
%   Fuß links:   © Can Siebert · 2026 — Fuß rechts: Seite X von Y
%   Titelseite (Seite 1) bleibt ohne Kopf-/Fußzeile (\handouttitel setzt
%   \thispagestyle{empty}).
% ---------------------------------------------------------------------------
\usepackage{fancyhdr}
\newcommand{\bfwKopfzeileText}{}
\newcommand{\kopfzeile}[1]{\renewcommand{\bfwKopfzeileText}{#1}}
\pagestyle{fancy}
\fancyhf{}
\fancyhead[L]{\small\bfseries\color{bfwPrimaryDark}Wirtschaftsfachwirt}
\fancyhead[R]{\small\color{bfwInkSoft}\bfwKopfzeileText}
\renewcommand{\headrulewidth}{0.6pt}
\renewcommand{\headrule}{{\color{bfwPrimary}\hrule width\headwidth height 0.6pt}}
\fancyfoot[L]{\small\color{bfwInkSoft}\copyright\ Can Siebert\,\textperiodcentered\,2026}
\fancyfoot[R]{\small\color{bfwInkSoft}Seite \thepage\ von \pageref*{LastPage}}
% Auch \thispagestyle{plain} (z.B. durch \maketitle) soll leer bleiben:
\fancypagestyle{plain}{\fancyhf{}\renewcommand{\headrulewidth}{0pt}\renewcommand{\headrule}{}}

% ---------------------------------------------------------------------------
% \handouttitel{Obertitel}{Untertitel}{Kurs-Label}{Termin-Zeile}
%   Einheitlicher professioneller Titelblock: Dachzeile, farbige Headline,
%   Untertitel, farbige Trennlinie und dezente Meta-Zeile (Kurs/Termin/Dozent).
%   Ersetzt \title/\maketitle. Direkt nach \begin{document} aufrufen.
% ---------------------------------------------------------------------------
\newcommand{\handouttitel}[4]{%
  \thispagestyle{empty}%
  \begingroup
  \setlength{\parindent}{0pt}%
  {\bfwlabelfont\color{bfwPrimary}WIRTSCHAFTSFACHWIRT (IHK)\par}
  \vspace{0.55em}
  {\huge\bfseries\color{bfwPrimaryDark}#1\par}
  \vspace{0.5em}
  {\large\color{bfwInkSoft}#2\par}
  \vspace{0.9em}
  {\color{bfwPrimary}\rule{\linewidth}{2pt}}\par
  \vspace{0.55em}
  {\small\color{bfwInkSoft}#3\ \,\textperiodcentered\,\ #4\ \,\textperiodcentered\,\ Dozent: Can Siebert\par}
  \endgroup
  \vspace{1.4em}%
}


\kopfzeile{Aufgabenheft VWL Lektion 3}

\begin{document}
\handouttitel{Aufgabenheft Lektion~3}%
  {Konjunktur, Wirtschaftspolitik \& Außenwirtschaft --- Übungen im IHK-Klausurformat}%
  {1.1 Volkswirtschaftliche Grundlagen}%
  {Selbstlernmaterial zur Lektion vom Sa, 29.08.2026}


\begin{merksatz}[title={\bfseries So nutzen Sie dieses Aufgabenheft}]
Bearbeiten Sie alle Aufgaben zunächst \emph{ohne} in die Lösungen zu schauen. Schreiben
Sie Ihre Antworten in vollständigen Sätzen --- die IHK-Prüfung verlangt formulierte
Begründungen. Beachten Sie die \emph{Operatoren} (nennen, erläutern, darstellen,
beurteilen) --- sie geben den geforderten Antwort-Tiefegrad vor. Erst danach öffnen Sie
den Lösungsteil ab Seite~\pageref{loesungen} und vergleichen. Dieses Heft deckt zugleich
zentralen Prüfungsstoff des gesamten Blocks~1.1 ab. Ergänzend: Übungsband Volkswirtschaft,
Kapitel~3 und~4.
\end{merksatz}

\tableofcontents
\vspace{1em}
\hrule

\section{Aufgaben}

\subsection*{Aufgabe~1 --- Konjunkturzyklus und Indikatoren (10 Punkte)}

Jede volkswirtschaftliche Entwicklung ist Schwankungen unterworfen.

\begin{enumerate}[label=(\alph*)]
  \item \textbf{Skizzieren} Sie einen Konjunkturzyklus grafisch und \textbf{beschriften}
    Sie die vier Konjunkturphasen. \hfill (4~P.)
  \item \textbf{Unterscheiden} Sie Konjunkturschwankung und Saisonschwankung. \hfill (3~P.)
  \item \textbf{Nennen} Sie drei Präsenzindikatoren, die zur Beschreibung der aktuellen
    Konjunkturlage herangezogen werden können. \hfill (3~P.)
\end{enumerate}

\subsection*{Aufgabe~2 --- Ziele und Instrumente der EZB (12 Punkte)}

Die Europäische Zentralbank (EZB) ist für die Preisniveaustabilität in der Eurozone
verantwortlich. Um diese zu erreichen, legt sie die Leitzinsen fest und bedient sich
weiterer geldpolitischer Instrumente.

\begin{enumerate}[label=(\alph*)]
  \item \textbf{Nennen} Sie das vorrangige Ziel der EZB und \textbf{erklären} Sie, wann
    dieses Ziel als erreicht gilt. \hfill (3~P.)
  \item \textbf{Geben} Sie die drei Leitzinsen der EZB \textbf{an}. \hfill (3~P.)
  \item \textbf{Beschreiben} Sie das Hauptrefinanzierungsgeschäft (Rhythmus, Laufzeit,
    Technik, Bedeutung). \hfill (4~P.)
  \item \textbf{Erklären} Sie den Begriff Mindestreserve. \hfill (2~P.)
\end{enumerate}

\subsection*{Aufgabe~3 --- EZB-Zinserhöhung und ihre Wirkungen (12 Punkte)}

Die Inflationsrate im Euroraum liegt deutlich über dem Zielwert von knapp zwei Prozent.
Der EZB-Rat beschließt daraufhin eine kräftige Erhöhung der Leitzinsen.

\begin{enumerate}[label=(\alph*)]
  \item \textbf{Stellen} Sie die Wirkungskette dieser Zinserhöhung bis zum Preisniveau
    \textbf{dar} (mindestens fünf Stufen). \hfill (6~P.)
  \item \textbf{Erläutern} Sie zwei mögliche negative Nebenwirkungen der Zinserhöhung auf
    Wachstum und Beschäftigung. \hfill (4~P.)
  \item \textbf{Nennen} Sie zwei Grenzen (Störfaktoren), an die die Geldpolitik der EZB
    stoßen kann. \hfill (2~P.)
\end{enumerate}

\subsection*{Aufgabe~4 --- Antizyklische Fiskalpolitik (12 Punkte)}

\begin{enumerate}[label=(\alph*)]
  \item \textbf{Erklären} Sie das Grundprinzip der antizyklischen Fiskalpolitik nach
    John Maynard Keynes und nennen Sie die gesetzliche Grundlage in Deutschland. \hfill (4~P.)
  \item \textbf{Nennen} Sie je zwei fiskalpolitische Maßnahmen für die Boomphase und für
    die Rezession. \hfill (4~P.)
  \item \textbf{Erklären} Sie den Begriff „automatische Stabilisatoren`` anhand zweier
    Beispiele. \hfill (4~P.)
\end{enumerate}

\subsection*{Aufgabe~5 --- Tarifpolitik (10 Punkte)}

\begin{enumerate}[label=(\alph*)]
  \item \textbf{Erklären} Sie den Begriff Tarifautonomie und \textbf{nennen} Sie die
    Rechtsgrundlage. \hfill (3~P.)
  \item \textbf{Unterscheiden} Sie Lohn- bzw. Gehaltstarifvertrag und Manteltarifvertrag
    nach ihren Regelungsinhalten. \hfill (4~P.)
  \item \textbf{Stellen} Sie die Lohnargumentation der Gewerkschaften (Kaufkrafttheorie)
    der produktivitätsorientierten Lohnpolitik der Arbeitgeber \textbf{gegenüber}. \hfill (3~P.)
\end{enumerate}

\subsection*{Aufgabe~6 --- Arbeitsmarktpolitik (8 Punkte)}

\begin{enumerate}[label=(\alph*)]
  \item \textbf{Unterscheiden} Sie aktive und passive Arbeitsmarktpolitik nach ihrer
    Zielsetzung. \hfill (4~P.)
  \item \textbf{Nennen} Sie je zwei Instrumente der aktiven und der passiven
    Arbeitsmarktpolitik. \hfill (4~P.)
\end{enumerate}

\subsection*{Aufgabe~7 --- Umweltpolitik (10 Punkte)}

\begin{enumerate}[label=(\alph*)]
  \item \textbf{Erklären} Sie Verursacherprinzip und Gemeinlastprinzip mit je einem
    Beispiel. \hfill (4~P.)
  \item \textbf{Erklären} Sie die Funktionsweise des Zertifikatehandels am Beispiel des
    EU-Emissionshandels. \hfill (4~P.)
  \item \textbf{Nennen} Sie zwei Beispiele für umweltpolitische Auflagen (Ordnungsrecht).
    \hfill (2~P.)
\end{enumerate}

\subsection*{Aufgabe~8 --- Transfer: Nachfrage- vs. angebotsorientierte Wirtschaftspolitik (14 Punkte)}

Deutschland befindet sich in einer Rezession: Das reale BIP schrumpft um 1,5~\%, die
Arbeitslosigkeit steigt. Im Bundestag stehen sich zwei Positionen gegenüber.
\emph{Partei~A} fordert ein kreditfinanziertes staatliches Investitionsprogramm für
Infrastruktur und höhere Transferleistungen für Haushalte. \emph{Partei~B} fordert
Steuersenkungen für Unternehmen, den Abbau von Bürokratie und Regulierungen sowie moderate
Lohnabschlüsse.

\begin{enumerate}[label=(\alph*)]
  \item \textbf{Ordnen} Sie beide Positionen der jeweiligen wirtschaftspolitischen
    Konzeption \textbf{zu} und \textbf{nennen} Sie deren theoretische Begründer bzw.
    prominente Vertreter. \hfill (4~P.)
  \item \textbf{Erläutern} Sie die Wirkungslogik beider Konzepte (Ziel, Ansatzpunkt, Rolle
    des Staates). \hfill (6~P.)
  \item \textbf{Beurteilen} Sie beide Konzepte, indem Sie je zwei Schwächen (Kritikpunkte)
    darlegen. \hfill (4~P.)
\end{enumerate}

\subsection*{Aufgabe~9 --- Außenwirtschaft und EU (14 Punkte)}

\begin{enumerate}[label=(\alph*)]
  \item \textbf{Unterscheiden} Sie tarifäre und nichttarifäre Handelshemmnisse und
    \textbf{nennen} Sie je zwei Beispiele. \hfill (4~P.)
  \item \textbf{Nennen} Sie zwei Vorteile des Freihandels und ein Motiv für
    protektionistische Maßnahmen. \hfill (3~P.)
  \item \textbf{Beschreiben} Sie die vier Grundfreiheiten des Europäischen Binnenmarktes
    jeweils anhand eines Beispiels. \hfill (4~P.)
  \item \textbf{Nennen} Sie drei Konvergenzkriterien für den Beitritt zur Eurozone (mit
    Referenzwerten). \hfill (3~P.)
\end{enumerate}

\newpage
\section{Lösungen}\label{loesungen}

\begin{loesung}[title={Lösung Aufgabe~1 --- Konjunkturzyklus und Indikatoren}]
\textbf{(a)} Die Skizze zeigt eine Wellenlinie (gesamtwirtschaftliche Aktivität, z.\,B.
reales BIP) um einen steigenden Wachstumstrend mit den vier Phasen:
(1)~\emph{Aufschwung (Expansion)} --- steigende Nachfrage, Produktion, Beschäftigung;
(2)~\emph{Hochkonjunktur (Boom)} --- volle Kapazitätsauslastung, steigende Preise und
Zinsen; (3)~\emph{Abschwung (Rezession)} --- sinkende Nachfrage, Auftragsrückgänge,
steigende Arbeitslosigkeit; (4)~\emph{Tiefphase (Depression)} --- geringe Auslastung,
hohe Arbeitslosigkeit, niedrige Preise und Zinsen.

\textbf{(b)} \emph{Konjunkturschwankungen} sind mehrjährige, wellenförmige Schwankungen
der gesamtwirtschaftlichen Aktivität um den langfristigen Trend.
\emph{Saisonschwankungen} sind kurzfristige, regelmäßig wiederkehrende, jahreszeitlich
bedingte Schwankungen (z.\,B. Bauwirtschaft im Winter, Tourismus im Sommer).

\textbf{(c)} Z.\,B. reales Bruttoinlandsprodukt, Kapazitätsauslastung,
Industrieproduktion, Zahl der offenen Stellen, Kurzarbeit (drei Nennungen genügen).
\end{loesung}

\begin{loesung}[title={Lösung Aufgabe~2 --- Ziele und Instrumente der EZB}]
\textbf{(a)} Vorrangiges Ziel ist die \emph{Preisniveaustabilität} in der Eurozone. Sie
gilt als erreicht, wenn der Anstieg des harmonisierten Verbraucherpreisindex (HVPI) unter,
aber nahe zwei Prozent liegt --- damit sind sowohl Inflation als auch Deflation
ausgeschlossen.

\textbf{(b)} (1)~Zinssatz für die Hauptrefinanzierungsgeschäfte (wichtigster Leitzins),
(2)~Zinssatz der Spitzenrefinanzierungsfazilität (Übernachtkredit, Obergrenze des
Zinskorridors), (3)~Zinssatz der Einlagefazilität (Übernachtanlage, Untergrenze).

\textbf{(c)} Das Hauptrefinanzierungsgeschäft (Haupttender) ist eine befristete
Transaktion, die im \emph{wöchentlichen} Rhythmus mit einer Laufzeit von \emph{einer
Woche} angeboten wird. Die EZB kauft bzw. verkauft Wertpapiere auf Basis einer
Rückkaufvereinbarung (Pensionsgeschäfte) oder gewährt Kredite gegen Verpfändung von
Wertpapieren. Die Banken decken hierüber den größten Teil ihres Refinanzierungsbedarfs;
der festgelegte Zinssatz ist der wichtigste Leitzins im Euroraum und beeinflusst sämtliche
Zinsen der Banken.

\textbf{(d)} Die Mindestreserve verpflichtet die Geschäftsbanken, bestimmte prozentuale
Anteile der bei ihnen angelegten Sicht-, Termin- und Spareinlagen als Pflichtguthaben auf
Konten bei der Zentralbank zu unterhalten.
\end{loesung}

\begin{loesung}[title={Lösung Aufgabe~3 --- EZB-Zinserhöhung und ihre Wirkungen}]
\textbf{(a)} Wirkungskette (restriktive Geldpolitik): Leitzinserhöhung $\Rightarrow$
(1)~Refinanzierung bei der EZB wird für die Geschäftsbanken teurer $\Rightarrow$
(2)~die Banken geben die höheren Zinsen an Unternehmen und Haushalte weiter, das
Zinsniveau steigt (zugleich steigt die Sparneigung) $\Rightarrow$ (3)~die Kreditnachfrage
sinkt, die Kreditmöglichkeiten werden eingeschränkt $\Rightarrow$ (4)~kreditfinanzierte
Investitionen der Unternehmen und Konsumentenkredite der Haushalte gehen zurück
$\Rightarrow$ (5)~die gesamtwirtschaftliche Nachfrage sinkt $\Rightarrow$ (6)~der
Preisanstieg wird verringert, die Inflation bekämpft (Konjunkturdämpfung).

\textbf{(b)} (1)~\emph{Wachstum:} Unterbleiben Investitionen, sinkt die
gesamtwirtschaftliche Nachfrage --- das Wirtschaftswachstum (BIP) geht zurück, im
Extremfall droht eine Rezession. (2)~\emph{Beschäftigung:} Bei sinkender Produktion bauen
Unternehmen Personal ab bzw. stellen nicht ein --- die Arbeitslosigkeit steigt. (Auch
denkbar: höhere Finanzierungskosten für den Staat, Belastung von Bauwirtschaft und
Immobilienmarkt.)

\textbf{(c)} Z.\,B.: Internationale Finanzströme können trotz restriktiver Politik
Inflation erzeugen; hohe Lohnabschlüsse über dem Produktivitätsfortschritt wirken
preistreibend; unklar, ob/wann Banken Zinsänderungen weitergeben; erhöht der Staat
gleichzeitig schuldenfinanziert seine Ausgaben, verpufft die restriktive Wirkung (zwei
Nennungen genügen).
\end{loesung}

\begin{loesung}[title={Lösung Aufgabe~4 --- Antizyklische Fiskalpolitik}]
\textbf{(a)} Der Staat gestaltet seine Einnahmen- und Ausgabenpolitik so, dass sie dem
Konjunkturverlauf \emph{entgegenwirkt}: Im Boom dämpft er (restriktiv), in der Rezession
stützt er die Nachfrage (expansiv). Konjunkturschwankungen sollen so ausgeglichen werden.
Gesetzliche Grundlage ist das \emph{Stabilitätsgesetz (StabG)} von 1967; theoretische
Grundlage die Lehre von John Maynard Keynes.

\textbf{(b)} \emph{Boom (restriktiv):} Staatsausgaben senken (Investitionen, Subventionen,
Transferleistungen zurückfahren), Einnahmen erhöhen (Steuererhöhungen), Budgetüberschüsse
als Rücklagen bilden. \emph{Rezession (expansiv):} Ausgaben erhöhen (öffentliche
Investitions- und Konjunkturprogramme), Einnahmen senken (Steuererleichterungen) ---
notfalls kreditfinanziert (Deficit Spending). Je zwei Nennungen genügen.

\textbf{(c)} Automatische (eingebaute) Stabilisatoren wirken antizyklisch, \emph{ohne}
dass eine politische Entscheidung nötig ist: (1)~Die \emph{progressive Einkommensteuer}
entzieht im Boom überproportional Kaufkraft und entlastet in der Rezession.
(2)~\emph{Arbeitslosengeld/Sozialleistungen} stützen in der Rezession die Nachfrage der
Haushalte, während im Boom Beitragszahlungen die Kaufkraft dämpfen.
\end{loesung}

\begin{loesung}[title={Lösung Aufgabe~5 --- Tarifpolitik}]
\textbf{(a)} Tarifautonomie bedeutet, dass die Tarifpartner (Gewerkschaften,
Arbeitgeberverbände bzw. einzelne Arbeitgeber) Tarifverträge frei von staatlichen
Weisungen aushandeln und abschließen; der Staat ist zur Neutralität verpflichtet.
Rechtsgrundlage: Art.~9 Abs.~3 GG (Koalitionsfreiheit).

\textbf{(b)} \emph{Lohn- und Gehalts(rahmen)tarifverträge} regeln die Entgeltfindung:
Höhe des Entgelts, Eingruppierung, Zuschläge und Zulagen; sie haben meist kurze
Laufzeiten. \emph{Manteltarifverträge} regeln die arbeitsrechtlichen Rahmenbedingungen:
Verteilung der Wochenarbeitszeit, Pausen, Urlaub, Schichtarbeit, Freistellungen,
Altersversorgung; sie gelten meist langfristig.

\textbf{(c)} \emph{Gewerkschaften (Kaufkrafttheorie, nachfrageorientiert):} Höhere Löhne
steigern die Massenkaufkraft und damit die Nachfrage; das regt Produktion und
Investitionen an und verbessert die Beschäftigung. \emph{Arbeitgeber
(produktivitätsorientiert, angebotsorientiert):} Löhne sind ein Kostenfaktor;
Lohnsteigerungen sollen den Produktivitätsfortschritt nicht übersteigen, um
Lohnkosteninflation und lohnkostenbedingte Arbeitslosigkeit (Rationalisierung,
Verlagerung) zu vermeiden.
\end{loesung}

\begin{loesung}[title={Lösung Aufgabe~6 --- Arbeitsmarktpolitik}]
\textbf{(a)} Die \emph{passive} Arbeitsmarktpolitik mildert die wirtschaftlichen Folgen
bereits eingetretener oder unmittelbar drohender Arbeitslosigkeit bzw. geringer
Beschäftigung (Absicherung des Einkommens). Die \emph{aktive} Arbeitsmarktpolitik setzt
vorher an: Sie soll Erwerbspersonen in die Lage versetzen, im Arbeitsprozess zu bleiben
oder schnellstmöglich und dauerhaft wieder einzutreten (Vorbeugung und Eingliederung).

\textbf{(b)} \emph{Passiv:} Arbeitslosengeld I und II (Bürgergeld), Kurzarbeitergeld,
Insolvenzgeld. \emph{Aktiv:} Beratung und Vermittlung, Trainings- und
Qualifizierungsmaßnahmen, Förderung der Aus- und Weiterbildung, Umschulung,
Existenzgründungsförderung, Mobilitätshilfen. Je zwei Nennungen genügen.
\end{loesung}

\begin{loesung}[title={Lösung Aufgabe~7 --- Umweltpolitik}]
\textbf{(a)} \emph{Verursacherprinzip:} Die Kosten für Vermeidung, Beseitigung oder
Ausgleich von Umweltbelastungen trägt derjenige, der sie verursacht --- externe Kosten
werden internalisiert (Beispiel: Abwasserabgabe, Rücknahmepflicht für Umverpackungen).
\emph{Gemeinlastprinzip:} Ist der Verursacher nicht eindeutig feststellbar oder gehen
Schäden zulasten der Allgemeinheit, übernimmt der Staat die Maßnahmen und finanziert sie
über Steuern (Beispiel: kommunale Kläranlagen, Mülldeponien).

\textbf{(b)} Unternehmen dürfen Umweltbelastungen nur gegen eine Berechtigung in Form von
\emph{Umweltzertifikaten} verursachen. Ein Zertifikat gestattet eine bestimmte Menge
Schadstoffausstoß pro Jahr. Die Gesamtmenge ist begrenzt (Cap); die Emissionsrechte werden
börsenmäßig gehandelt bzw. versteigert. Wer wenig emittiert, kann Zertifikate verkaufen ---
so entsteht ein marktwirtschaftlicher Anreiz, Emissionen dort zu vermeiden, wo es am
günstigsten ist (EU-Emissionshandel für CO\textsubscript{2}).

\textbf{(c)} Z.\,B. Grenzwerte für den Schadstoffausstoß, FCKW-Verbot, Umweltplaketten
für Kraftfahrzeuge beim Einfahren in Umweltzonen (zwei Nennungen genügen).
\end{loesung}

\begin{loesung}[title={Lösung Aufgabe~8 --- Nachfrage- vs. angebotsorientierte Wirtschaftspolitik}]
\textbf{(a)} \emph{Partei~A} vertritt die \textbf{nachfrageorientierte} Wirtschaftspolitik
(keynesianische Nachfragesteuerung, Fiskalismus) --- Begründer: John Maynard Keynes
(1883--1946). \emph{Partei~B} vertritt die \textbf{angebotsorientierte} Wirtschaftspolitik
(Neoklassik/Angebotstheorie, monetaristisch geprägt) --- prominenter Vertreter: Milton
Friedman (1912--2006).

\textbf{(b)} \emph{Nachfrageorientiert:} Ursache der Krise ist eine unzureichende
gesamtwirtschaftliche Nachfrage. Ziel ist die kurzfristige Beseitigung der
Gleichgewichtsstörung (Symptombekämpfung). Der Staat übernimmt eine \emph{aktive} Rolle:
Er stärkt die Massenkaufkraft (Transfers, Lohnerhöhungen), erhöht den Staatskonsum über
öffentliche Ausgabenprogramme und verschuldet sich dafür notfalls (Deficit Spending) ---
antizyklische Politik. \emph{Angebotsorientiert:} Wachstum und Beschäftigung hängen von
den Angebotsbedingungen ab; Motor sind die Investitionen der Unternehmen. Ziel ist die
mittel- bis langfristige Beseitigung gleichgewichtsstörender Auslösefaktoren
(Ursachenbekämpfung). Der Staat verhält sich \emph{passiv} und verbessert nur die
Rahmenbedingungen: Steuersenkungen, Deregulierung, Kostendämpfung, Abbau von Subventionen
--- bessere Renditeerwartungen führen zu Investitionen und neuen Arbeitsplätzen.

\textbf{(c)} \emph{Kritik an A (nachfrageorientiert):} Timelags (Wirkung kommt oft zu
spät), Crowding-out (Verdrängung privater Investitionen), steigende Staatsquote und
wachsende Staatsverschuldung. \emph{Kritik an B (angebotsorientiert):} einseitige
Betonung der Kosten --- Löhne sind auch ein Nachfragefaktor; höhere Gewinne erhöhen zwar
die Investitionsfähigkeit, aber nicht automatisch die Investitionsbereitschaft (bei
schlechter Absatzerwartung wird nicht investiert). Je zwei Punkte genügen.
\end{loesung}

\begin{loesung}[title={Lösung Aufgabe~9 --- Außenwirtschaft und EU}]
\textbf{(a)} \emph{Tarifäre Handelshemmnisse} wirken über den \emph{Preis} der Ein- und
Ausfuhrgüter (preispolitische Instrumente): z.\,B. Importzölle, Exportsubventionen,
Steuertarife. \emph{Nichttarifäre Handelshemmnisse} haben keine fiskalische Funktion,
sondern wirken mengenmäßig oder administrativ: z.\,B. Importkontingente, Ein- und
Ausfuhrverbote, Embargos, technische Normen und Genehmigungsverfahren.

\textbf{(b)} \emph{Vorteile des Freihandels:} vielfältiges Güterangebot für Verbraucher
(Qualitäten, Preise), Spezialisierungsvorteile und Kostensenkungen, Absatzmärkte im In-
und Ausland, Innovationsförderung durch internationalen Wettbewerb (zwei Nennungen).
\emph{Motiv für Protektionismus:} Schutz der einheimischen Wirtschaft vor Billigimporten
bzw. vor konjunkturellen/strukturellen Nachteilen, Vermeidung von Monostrukturen und
Auslandsabhängigkeit (eine Nennung).

\textbf{(c)} (1)~\emph{Freier Personenverkehr:} z.\,B. freie Arbeitsplatzwahl und
Niederlassungsfreiheit in der ganzen EU, Abbau der Grenzkontrollen (Schengen).
(2)~\emph{Freier Warenverkehr:} keine Zölle oder mengenmäßigen Beschränkungen im
EU-Handel, harmonisierte technische Normen. (3)~\emph{Freier Dienstleistungsverkehr:}
Selbstständige und freie Berufe dürfen ihre Leistungen EU-weit anbieten (z.\,B. ein
deutscher Handwerksbetrieb in Österreich). (4)~\emph{Freier Geld- und Kapitalverkehr:}
Öffnung der Finanzmärkte, erleichterter Zahlungsverkehr (SEPA), freie Kapitalanlage im
EU-Ausland.

\textbf{(d)} Z.\,B.: Inflationsrate höchstens 1,5~Prozentpunkte über dem Mittelwert der
drei preisstabilsten EU-Staaten; Nettoneuverschuldung nicht dauerhaft über 3~\% des BIP;
Gesamtschuldenstand i.\,d.\,R. höchstens 60~\% des BIP; langfristiger Zinssatz höchstens
2~Prozentpunkte über dem Mittelwert der drei preisstabilsten Staaten; zwei Jahre stabiler
Wechselkurs im WKM~II (drei Nennungen genügen).
\end{loesung}

\vspace{1em}
\hrule
\vspace{0.8em}
\noindent\textsf{\small Weiterüben: Übungsband Volkswirtschaft Kap.~3 und~4 (mit Lösungsteil), Quiz und
Karteikarten der Lektionen 1--3. Nächster Termin: Start des BWL-Teils (Betriebliche Funktionen~I, Sa 05.09.2026).}

\end{document}
